\section{Problem Formulation}

\subsection{Baseline Problem: Convergence Optimization}

The fundamental objective of federated learning is to minimize the global loss function through iterative client selection and model aggregation. We first formulate the baseline optimization problem:

\textbf{Problem P1 (Convergence-Only):}
\begin{equation}
\min_{\{a_n(t)\}} \quad F(w_T)
\end{equation}
\begin{subequations}
\begin{align}
\text{s.t.} \quad & \sum_{n=1}^{N} a_n(t) = K, \quad \forall t \in \{1, \ldots, T\}, \\
& a_n(t) \in \{0, 1\}, \quad \forall n, t,
\end{align}
\end{subequations}

where $a_n(t)$ is the binary selection indicator for client $n$ at round $t$, and $w_T$ is the global model after $T$ rounds.

\textbf{Limitation:} Problem P1 ignores client heterogeneity and resource constraints. Random selection (FedAvg) treats all clients equally, leading to suboptimal convergence in non-IID settings and potential energy depletion for frequently selected clients.

\subsection{Contribution-Aware Problem with Energy Constraints}

To address client heterogeneity, we introduce Shapley values to quantify marginal contributions and add energy sustainability constraints:

\textbf{Problem P2 (Shapley + Energy):}
\begin{subequations}
\begin{align}
\max_{\{a_n(t)\}} \quad & \sum_{t=1}^{T} \sum_{n=1}^{N} a_n(t) \cdot \varphi_n(t) \\
\text{s.t.} \quad & \sum_{n=1}^{N} a_n(t) = K, \quad \forall t, \\
& \frac{1}{T} \sum_{t=1}^{T} a_n(t) E_n(t) \leq \bar{E}_n, \quad \forall n, \\
& a_n(t) \in \{0, 1\},
\end{align}
\end{subequations}

where $\varphi_n(t)$ is the Shapley value of client $n$ at round $t$, $E_n(t) = \sigma^2 / |h_n(t)|^2$ is the transmission energy consumption, and $\bar{E}_n$ is the average energy budget.

\textbf{Shapley Value Definition:} For client $n$, the Shapley value quantifies its marginal contribution:
\begin{equation}
\varphi_n(t) = \sum_{S \subseteq \mathcal{N} \setminus \{n\}} \frac{|S|! (N - |S| - 1)!}{N!} [v(S \cup \{n\}) - v(S)]
\end{equation}
where $v(S)$ is the utility function defined as the negative validation loss of the aggregated model from coalition $S$.

\textbf{Limitation:} Problem P2 couples decisions across all rounds through the average energy constraint, making it computationally intractable for online decision-making.

\subsection{Lyapunov-Based Online Optimization}

To enable online decision-making, we reformulate P2 using Lyapunov optimization, which decouples the long-term energy constraint into per-round decisions.

\textbf{Virtual Queue Construction:} For each client $n$, define a virtual energy queue that tracks cumulative energy constraint violations:
\begin{equation}
Q_n(t+1) = \max\{Q_n(t) + E_n(t) - \bar{E}, \; 0\}
\end{equation}
where $E_n(t)$ is the actual energy consumed by client $n$ in round $t$ (zero if not selected), and $\bar{E}$ is the per-round energy budget. The queue increases when a client's energy consumption exceeds the budget, signaling that the client has been consuming excessive energy.

\textbf{Lyapunov Function:}
\begin{equation}
L(Q(t)) = \frac{1}{2} \sum_{n=1}^{N} Q_n^2(t)
\end{equation}

\textbf{Problem P3 (Per-Round Optimization):} At each round $t$, minimize the drift-plus-penalty:
\begin{subequations}
\begin{align}
\min_{\{a_n(t)\}} \quad & \Delta(Q(t)) + V \cdot \mathbb{E}[-\text{Utility}(t)] \\
\text{s.t.} \quad & \sum_{n=1}^{N} a_n(t) = K, \\
& a_n(t) \in \{0, 1\},
\end{align}
\end{subequations}

where $\Delta(Q(t)) = \mathbb{E}[L(Q(t+1)) - L(Q(t))]$ is the Lyapunov drift, and $V > 0$ controls the tradeoff between queue stability (energy) and utility (accuracy).

\subsection{Practical Solution: Composite Scoring}

To solve P3 efficiently, we design a composite scoring function derived from the drift-plus-penalty framework. By upper-bounding the Lyapunov drift term, the per-round optimization reduces to maximizing the following score for each client:

\textbf{Composite Score:}
\begin{equation}
\text{Score}_n(t) = V \cdot \hat{\varphi}_n(t) - \hat{Q}_n(t) \cdot \hat{E}_n(t)
\end{equation}
where $\hat{\varphi}_n(t) = (\varphi_n(t) - \varphi_{\min}(t)) / (\varphi_{\max}(t) - \varphi_{\min}(t))$ is the min-max normalized Shapley value, $\hat{E}_n(t) = E_n(t) / E_{\max}(t)$ is the normalized energy consumption, and $\hat{Q}_n(t) = Q_n(t) / \max_{i} Q_i(t)$ is the normalized per-client virtual queue. Normalizing the queue prevents its magnitude from overwhelming the Shapley term as energy violations accumulate over time. The parameter $V > 0$ controls the tradeoff: larger $V$ prioritizes Shapley-based client quality, while the normalized queue $\hat{Q}_n(t)$ automatically penalizes clients that have accumulated excessive energy consumption.

\textbf{Greedy Selection:} At each round, select the top-$K$ clients with highest composite scores among energy-available clients.

\textbf{Secure Transmission:} Encrypt gradient updates using AES-256-GCM before uploading to the server; the server decrypts ephemerally for Shapley computation and aggregation, then immediately destroys plaintext.

